% =====================================================================
% paper/sections/length_bias_react_shaping.tex
%
% PROVENANCE. Berkeley B-mining ledger row 13 (iter145, 2026-07-04):
% ReAct-style intermediate-credit-shaping diagnostic on bfclv4
% tool-use (F24 L2, Shunyu Yao; arXiv:2210.03629).
% Driver: scripts/berkeley/react_reasoning_act.py
% Inputs:  experiments/results/bfclv4_tool_use.tsv
% Outputs: experiments/results/berkeley/react_{dense_vs_sparse_step,
%          intervention_gap,zvf_reduction}.tsv
%          experiments/results/berkeley/react_summary.json
%
% Ledger verdict: SUGGESTIVE (2/4 pre-registered hypotheses DECISIVE:
% H1 Cohen's d = 0.667, H4 ZVF drop = 0.10; H2 and H3 NULL).
% Framed accordingly: a suggestive diagnostic, not a validated result.
% =====================================================================

\subsection{A Suggestive Tool-Use Diagnostic: ReAct-Style Intermediate-Credit Shaping}
\label{sec:length-bias-react-shaping}

ReAct~\citep{yao2023react} interleaves reasoning traces with
task actions and attributes its gains to intermediate states
becoming credit-bearing rather than rolling up into a single
end-of-trajectory signal. Transcribed to RL post-training, the
analogous intervention is \emph{intermediate-credit (dense)
reward shaping}---directly relevant to this pillar, where
end-anchored sparse credit is the hypothesised driver of both
the length blow-out and the zero-signal failure modes. We test
the transcription on the BFCLv4 tool-use rollout summary
(Qwen~0.5B; 2~seeds, 10~step-level observations;
\texttt{experiments/results/berkeley/react\_dense\_vs\_sparse\_step.tsv}),
where the sparse- and dense-reward columns differ only in reward
sparsification.

\paragraph{Pre-registered hypotheses.}
Four reward-summary-level hypotheses, with thresholds fixed
before measurement: \textbf{H1}~dense beats sparse on average
rollout reward (decisive iff Cohen's $d > 0.4$);
\textbf{H2}~the first-half vs last-half reward delta is larger
under dense (credit is used earlier); \textbf{H3}~the fraction
of seed-trajectories with identically-zero reward drops under
dense (decisive iff drop $> 0.05$); \textbf{H4}~the
zero-variance fraction (ZVF; the Pillar-2 signal-availability
variable) drops under dense (decisive iff drop $> 0.05$).

\paragraph{Result: 2/4 decisive.}
H1 is \textbf{decisive}: dense lifts the mean of seed-mean
rewards from $0.1125$ to $0.1863$ ($\Delta = +0.074$, Cohen's
$d = 0.667$;
\texttt{experiments/results/berkeley/react\_intervention\_gap.tsv}).
H4 is \textbf{decisive}: pooled ZVF drops from $0.40$ (sparse)
to $0.30$ (dense), a reduction of $0.10$
(\texttt{experiments/results/berkeley/react\_zvf\_reduction.tsv}).
The other two hypotheses are \textbf{null} on this data: the H2
half-life delta comes out $-0.044$, i.e.\ sign-opposite to the
prediction, and H3 cannot move because no seed on this rollout
has an identically-zero reward trace (sparse and dense zero-floor
fractions are both $0.0$;
\texttt{experiments/results/berkeley/react\_summary.json}).

\paragraph{Scope and verdict.}
With 2/4 hypotheses decisive, the mining ledger classifies this
finding as \textbf{suggestive}, not decisive, and we present it
strictly as a diagnostic. The evidence base is small ($n{=}2$
seeds) and is a re-analysis of an existing rollout summary
rather than a live ReAct intervention with explicit
thought/action tokens; the pre-registered promotion criterion
($n \ge 5$ seeds on longer trajectories, flipping H2/H3) has not
been run. Within that scope, the two decisive results point the
same way as this pillar's mechanism story: dense intermediate
credit reduces signal starvation on tool-use---raising realised
reward while restoring within-group contrast (lower ZVF)---which
is what one expects if the length and zero-signal pathologies
documented here are downstream of end-anchored credit
assignment.
